\documentclass[11pt]{article}
\usepackage[a4paper,margin=1in]{geometry}
\usepackage{amsmath,amssymb}
\usepackage{tcolorbox}
\usepackage{hyperref}

\title{Backpropagation -- Simple Recall Notes}
\author{}
\date{}

\begin{document}
\maketitle

\section*{1. The Big Idea}

A neural network has two main jobs:

\begin{enumerate}
\item \textbf{Forward pass:} calculate the output.
\item \textbf{Backward pass:} find how the weights should change to reduce the error.
\end{enumerate}

\[
\boxed{\text{Input}\rightarrow\text{Hidden}\rightarrow\text{Output}
\rightarrow\text{Error}\rightarrow\text{Backpropagation}}
\]

\begin{tcolorbox}
\textbf{Easy memory rule:}\\
Forward = ``What did the network predict?''\\
Backward = ``How did the weights contribute to the error?''
\end{tcolorbox}

\section*{2. What is a Hidden Unit?}

A hidden unit is simply a neuron between the input and output layers.

\[
\text{Input}\rightarrow\boxed{\text{Hidden}}\rightarrow\text{Output}
\]

We know the desired output, but we do not directly specify what each hidden unit should output. The network learns useful hidden representations through training.

\section*{3. Our Small Network}

Consider:

\[
\boxed{3\text{ inputs}\rightarrow2\text{ hidden neurons}\rightarrow1\text{ output}}
\]

Write the input as a column vector:

\[
X=
\begin{bmatrix}
x_1\\x_2\\x_3
\end{bmatrix}
\qquad (3\times1)
\]

and the input-to-hidden weights as:

\[
W_1=
\begin{bmatrix}
w_{11}&w_{12}&w_{13}\\
w_{21}&w_{22}&w_{23}
\end{bmatrix}
\qquad (2\times3)
\]

Therefore:

\[
\boxed{(2\times3)(3\times1)=(2\times1)}
\]

The middle numbers must match.

\section*{4. Why Does $W_1X$ Work?}

Suppose

\[
X=
\begin{bmatrix}
1\\0\\1
\end{bmatrix}
\]

and

\[
W_1=
\begin{bmatrix}
0.5&-1&0.8\\
-0.2&0.3&0.4
\end{bmatrix}.
\]

First hidden neuron:

\[
0.5(1)+(-1)(0)+0.8(1)=1.3
\]

Second hidden neuron:

\[
-0.2(1)+0.3(0)+0.4(1)=0.2
\]

Thus

\[
W_1X=
\begin{bmatrix}
1.3\\0.2
\end{bmatrix}.
\]

\section*{5. Forward Pass}

Add bias and activation:

\[
Z_1=W_1X+b_1
\]

\[
H=f(Z_1)
\]

Then:

\[
Z_2=W_2H+b_2
\]

\[
Y=f(Z_2)
\]

So:

\[
\boxed{X\rightarrow H\rightarrow Y}
\]

is the forward pass.

\section*{6. Error}

A simple squared-error function is:

\[
\boxed{E=\frac12(d-Y)^2}
\]

where $d$ is the desired output and $Y$ is the network output.

Learning tries to make $E$ smaller.

\section*{7. What Does Backpropagation Need to Find?}

For every weight $w$, we want:

\[
\boxed{\frac{\partial E}{\partial w}}
\]

In simple English:

\begin{quote}
If I change this weight a little, how does the error change?
\end{quote}

Then update the weight:

\[
\boxed{
w_{\text{new}}
=
w_{\text{old}}
-
\eta\frac{\partial E}{\partial w}
}
\]

where $\eta$ is the learning rate.

\section*{8. Why the Chain Rule?}

The error does not directly depend on an early weight.

Instead:

\[
w\rightarrow h\rightarrow y\rightarrow E.
\]

So we use the chain rule:

\[
\boxed{
\frac{\partial E}{\partial w}
=
\frac{\partial E}{\partial y}
\frac{\partial y}{\partial h}
\frac{\partial h}{\partial w}
}
\]

For a longer network, we simply keep multiplying the derivatives along the path.

\begin{tcolorbox}
\textbf{Remember:}\\
Chain rule = follow the path backward and multiply the local derivatives.
\end{tcolorbox}

\section*{9. One-Neuron Example}

For one hidden neuron:

\[
w_1\rightarrow z_h\rightarrow h\rightarrow z_y\rightarrow y\rightarrow E
\]

Backward:

\[
E\rightarrow y\rightarrow z_y\rightarrow h\rightarrow z_h\rightarrow w_1
\]

Therefore:

\[
\frac{\partial E}{\partial w_1}
=
\frac{\partial E}{\partial y}
\frac{\partial y}{\partial z_y}
\frac{\partial z_y}{\partial h}
\frac{\partial h}{\partial z_h}
\frac{\partial z_h}{\partial w_1}.
\]

Do not memorize this long equation yet. Remember the path.

\section*{10. Two Hidden Neurons}

Let

\[
H=
\begin{bmatrix}
h_1\\h_2
\end{bmatrix}.
\]

The output layer receives both hidden activations.

The output error is propagated backward to both hidden neurons.

The simple idea is:

\[
\boxed{
\text{hidden error}
=
\text{output error}
\times
\text{connecting weight}
\times
\text{activation derivative}
}
\]

This is how the network assigns responsibility for the final error to earlier units.

\section*{11. The Important Matrix Gradient}

We already know:

\[
W_1:(2\times3)
\]

and

\[
X:(3\times1).
\]

Forward:

\[
\boxed{
W_1X=(2\times3)(3\times1)=(2\times1)
}
\]

During backpropagation, let the hidden error vector be:

\[
\delta_h=
\begin{bmatrix}
\delta_1\\
\delta_2
\end{bmatrix}
\qquad(2\times1)
\]

and:

\[
X^T=
\begin{bmatrix}
x_1&x_2&x_3
\end{bmatrix}
\qquad(1\times3).
\]

Then:

\[
\boxed{
\frac{\partial E}{\partial W_1}
=
\delta_hX^T
}
\]

because

\[
\boxed{
(2\times1)(1\times3)=(2\times3)
}
\]

The gradient has the same shape as $W_1$.

\section*{12. The Shape Trick}

\[
\boxed{
\text{Forward: }W_1X
}
\]

\[
(2\times3)(3\times1)=(2\times1)
\]

and

\[
\boxed{
\text{Backward: }\delta_hX^T
}
\]

\[
(2\times1)(1\times3)=(2\times3)
\]

Therefore:

\[
\boxed{\text{gradient shape}=\text{weight shape}}
\]

This lets us update the whole weight matrix:

\[
\boxed{
W_1^{new}
=
W_1-\eta\frac{\partial E}{\partial W_1}
}
\]

\section*{13. The Whole Idea in One Picture}

\[
\boxed{
\text{FORWARD:}\quad
X\rightarrow H\rightarrow Y\rightarrow E
}
\]

\[
\boxed{
\text{BACKWARD:}\quad
E\rightarrow Y\rightarrow H\rightarrow W
}
\]

\begin{tcolorbox}
\centering
\textbf{Predict $\rightarrow$ Measure Error $\rightarrow$ Send Error Back $\rightarrow$ Update Weights}
\end{tcolorbox}

\section*{14. What to Memorize}

Do not memorize all the long equations yet.

Remember only:

\begin{enumerate}
\item Forward = calculate prediction.
\item Error = compare prediction with target.
\item Chain rule = follow the effect backward.
\item Gradient = tells us how a weight affects error.
\item Update = subtract the gradient.
\end{enumerate}

The core equations are:

\[
Z_1=W_1X+b_1
\]

\[
H=f(Z_1)
\]

\[
Z_2=W_2H+b_2
\]

\[
Y=f(Z_2)
\]

\[
E=\frac12(d-Y)^2
\]

\[
W_{\text{new}}
=
W_{\text{old}}
-
\eta\frac{\partial E}{\partial W}
\]

\end{document}
